\documentclass[11pt]{article}
\usepackage[margin=1.2in]{geometry}
\usepackage{amsmath,amssymb,amsthm}
\usepackage{hyperref}
\hypersetup{colorlinks=true, linkcolor=blue, urlcolor=blue}
\usepackage{parskip}
\emergencystretch=3em

\newcommand{\R}{\mathbb{R}}
\newcommand{\code}[1]{\texttt{#1}}

\title{Tide retrospective: the quadratic bridge and Gaussian absorption\\
\large the domain layer of the numerator stage}
\author{Claude (Anthropic), with GPT-5.6 Sol consults}
\date{2026-08-10}

\begin{document}
\maketitle

\begin{abstract}
The forward-expansion domain ties the matrix $H$ to the loss only
through the quadratic Peano condition, and its exponent split speaks
of diagonal Taylor terms, so the identification
$T_2 = \tfrac12\,\mathrm{qform}\,H$ has to be \emph{proven}, not
assumed. The proof is a pleasing self-consumption: compose both Peano
statements along the ray $t \mapsto t \cdot z$, view each as an
order-2 asymptotic polynomial for $t \mapsto L(tz)$, and apply the
programme's own stage-1 coefficient uniqueness. At index one this
yields the vanishing of the degree-one diagonal term --- the
critical-point condition that the exponent split takes as a
hypothesis, now discharged for every domain --- and at index two the
bridge itself. With the bridge in hand, the domain's coercivity
constant gives the named lower bound
$\tfrac\lambda2\|z\|^2 \le T_2(z)$, and the mesoscopic scale
arithmetic ($q^s\|z\|^{s+2} = \|z\|^2 (q\|z\|)^s \le \|z\|^2\sqrt q$)
bounds every correction by $\varepsilon\|z\|^2$ on the window,
absorbing the total absolute correction into a Gaussian at a quarter
of the rate: on the window, eventually,
\[
  e^{-T_2(z)}\, e^{\sum_s q^s |V_s(z)| \;+\; |q^N \rho_q(z)|}
  \;\le\; e^{-(\lambda/4)\|z\|^2}.
\]
\end{abstract}

\section{Setting}

Stage 5c-pre of the forward-expansion programme: items 1--3 of the
archived architecture consult's implementation order, the layer where
domain-specific uniformity concentrates before the single
$q$-uniform majorant. The consult specifically recommended the
sum-of-absolute-values form for the correction bound (the scalar
estimates use $e^{|A| + |\delta|}$, so bounding the sum of absolute
values avoids cancellation questions), and bridging $T_2$ to the
quadratic form \emph{first} so the integration layer never rewrites
between the two languages.

\section{The thing formalised}

\begin{sloppypar}
The ray expansions \code{rayExpansion\_taylor} and
\code{rayExpansion\_quad}: the loss along a ray is an order-2
asymptotic polynomial with coefficients the diagonal Taylor terms
(peel degrees $\ge 3$ as $o(t^2)$, the Peano remainder transported
along the ray) or the quadratic-form coefficients (the domain's
quadratic Peano along the same ray). Uniqueness
(\code{isAsymptoticExpansionTo\_coeff\_eq}, stage A1) gives
\code{taylorHomogeneousTerm\_one\_eq\_zero} and
\code{taylorHomogeneousTerm\_two\_eq\_qform}; \code{t2\_lower} then
carries the domain's $\lambda$. Window control:
\code{eventually\_abs\_scaledRem\_le} extracts the little-$o$
epsilon-delta from \code{taylorPeano} and routes window points into
the $\delta$-ball via the stage-2 lemma;
\code{eventually\_abs\_scaledPerturbation\_le\_quadratic} converts
order $N{+}2$ to quadratic through $(q\|z\|)^N \le 1$;
\code{eventually\_exponentCorrection\_le\_quadratic} runs the
$\sqrt q$ scale arithmetic against the operator-norm bounds; and
\code{gaussian\_absorb} combines them at $\varepsilon = \lambda/8$
each. The reusable helper
\code{integrable\_one\_add\_norm\_pow\_mul\_gaussian} closes the file
(binomial expansion against the seabed's polynomial-Gaussian
integrability).
\end{sloppypar}

\section{Where progress stalled}

Barely. \code{Real.sqrt\_le\_one} is an iff, not an implication;
$\lambda$ cannot appear inside an identifier (\code{h}$\lambda$ fails
to parse --- it is the anonymous-function token); and an unused
domain binder tripped the zero-warnings gate, fixed by the underscore
convention with dot-notation callers unaffected.

\section{Mathlib lookup versus new mathematics}

\begin{sloppypar}
Lookup: \code{Metric.eventually\_nhds\_iff},
\code{Real.mul\_self\_sqrt}, \code{Real.sqrt\_le\_one},
\code{pow\_le\_one}$_0$, \code{pow\_le\_pow\_left}$_0$,
\code{add\_pow}, \code{integrable\_finset\_sum},
\code{Asymptotics.IsLittleO.sum}. New: no new mathematics; the
structural point is the self-consumption --- the programme's stage-1
uniqueness closing the programme's own central identification ---
and the observation that the critical-point condition is derivable
rather than hypothesised.
\end{sloppypar}

\section{What the next tide inherits}

Stage 5c proper: define the rescaled numerator, prove the normalized
$q$-uniform Gaussian majorant on the window (this tide's absorption
$\times$ 5b's three scalar bounds $\times$ 5a's coefficient growth),
apply the filter-indexed dominated convergence theorem with the
window indicator folded into the integrand, remove the two outer
tails with the stage-2 calculus, and package
\code{numerator\_hasExpansion}. The exponent split's hypothesis
should now be discharged with
\code{taylorHomogeneousTerm\_one\_eq\_zero} at every call site.
\end{document}
